\section{Interpolation, boundary averaging, and completion}
\label{assembly:section}

We continue with a finite popular collection and its adjoint sum
\(P=\sum_s\phi_s\).  Thus
\(N=(|E|/|F|)^{1/2}\), \(L/2\le\log N\le2L\), and
\(S\lesssim L^C|E|\).  All statements in this section are uniform in
the finite collection, its depth, and the dimension of \(\mathcal H\).
We write \(\vartheta\) for the translation of the angular grid.  The
set \(Z_0\) and the functions \(P,P_q,h_T,n\) do not depend on
\(\vartheta\), whereas the smaller set \(Z=Z_\vartheta\), the angular
masses, and the analytic sums do.  Until the boundary-averaging step,
we fix \(\vartheta\) and suppress it in the notation.

\subsection{From the path estimate to analytic interpolation}

\begin{proposition}[Vertex estimate]\label{assembly:vertices}
For every nonempty root \(v\), every shift \(\sigma\), and every
parameter with
\[
 \Re\zeta=\mu e_\sigma,
 \qquad \|\Im\zeta\|_2\le L^2,
\]
one has
\begin{equation}\label{eq:src-29}
 \big\|n_v^{\alpha/p-1}A_v(\zeta)\big\|_{L^p(Z;\mathcal H)}
       \le L^C S^{1/p}.
\end{equation}
The constant and exponent are independent of the shift and the grid
translation.
\end{proposition}

\begin{proof}
Recall the nonnegative amount
\[
 V_v=n_v+n_v^{\alpha-p}\|A_v\|^p.
\]
On the set where \(n_v\le M\), Equation~\eqref{eq:src-17} gives
\[
 V_v\le L^{Cp}(1+M)^{\alpha-1}n_v.
\]
Here and below the bound on the imaginary parts is absorbed into
\(L^{Cp}\).  On the remaining set, apply
Lemma~\ref{hier:path-expansion}.  At least half of the root amount is
carried by good paths, and the terminal amount of a path ending at
\(R\) is at most
\(L^{Cp}(1+M)^{\alpha-1}h_R\).  Consequently,
\begin{align*}
 \int_Z V_v
 &\le L^{Cp}(1+M)^{\alpha-1}\int n_v\\
 &\quad+2L^{Cp}(1+M)^{\alpha-1}
   \sum_{R:\,s_R\in v}
   \int_{Z\cap\{\text{good }R\}}
        h_R\prod_{\text{active path of }R}\beta_J.
\end{align*}
The root and its descendants contain only their eligible tops; sums
over tops here and subsequently are understood with that convention.
Theorem~\ref{res:path-product}, namely
Equation~\eqref{eq:src-22}, bounds each integral in the second line by
\(L^{Cp}(1+M)^C\int h_R\).  Moreover,
\[
 \sum_R\int h_R\lesssim L^C S,
\]
by the definition of the dilated decaying weights.  We conclude that
\[
 \int_Z V_v\le L^{Cp}(1+M)^C S.
\]
Since \(M\) is a fixed power of \(L\), taking a \(p\)-th root and
using \(n_v^{\alpha-p}\|A_v\|^p\le V_v\) proves the result.
\end{proof}

\begin{lemma}[Interpolation of the shift parameters]
\label{assembly:interpolation}
Equation~\eqref{eq:src-29} also holds at
\[
 \zeta^{\mathrm c}=(\mu/K,\ldots,\mu/K).
\]
\end{lemma}

\begin{proof}
Let
\[
 \mathfrak S_\mu
 =\{a\in[0,\infty)^K:\ \textstyle\sum_\sigma a_\sigma=\mu\},
 \qquad
 \mathcal F_v(\zeta,z)=n_v(z)^{\alpha/p-1}A_v(\zeta,z).
\]
Every nonempty mass is positive.  For each fixed \(z\), the defining
sum is a finite sum of exponentials of linear functions of \(\zeta\),
with real logarithms of positive ratios as coefficients.  In
particular it is entire in \(\zeta\).

We first give a uniform bound on the tube over
\(\mathfrak S_\mu\), without using the vertex estimate.  From
Equations~\eqref{eq:src-17} and \eqref{eq:src-12}, and from
\(\int n_v\le\int n\lesssim L^C S\), we have, for
\(a\in\mathfrak S_\mu\) and \(\eta\in\R^K\),
\begin{align}
 \|\mathcal F_v(a+i\eta)\|_{L^p(Z)}
 &\le L^C(1+\mu+\|\eta\|_2)
           \left(\int_Z n_v^\alpha\right)^{1/p}\notag\\
 &\le L^C(1+\|\eta\|_2)
          (NL^C)^{(\alpha-1)/p}S^{1/p}.
 \label{assembly:crude-analytic-bound}
\end{align}
The factor involving \(N\) is at most \(\exp(C\sqrt L)\), since
\(\log N\le2L\), \(p=L^{1/2}\), and \(2<\alpha<3\) for
sufficiently large \(L\).

Multiply the analytic family by
\[
 \mathcal D(\zeta)=\exp\left(\sum_\sigma\zeta_\sigma^2\right).
\]
On the tube under consideration,
\[
 |\mathcal D(a+i\eta)|
       =\exp(\|a\|_2^2-\|\eta\|_2^2)
       \le\exp(\mu^2-\|\eta\|_2^2).
\]
At a vertex, Proposition~\ref{assembly:vertices} controls
\(\mathcal D\mathcal F_v\) whenever \(\|\eta\|_2\le L^2\).
For \(\|\eta\|_2>L^2\),
Equation~\eqref{assembly:crude-analytic-bound} gives the same bound:
indeed
\[
 (1+\|\eta\|_2)
  \exp(\mu^2-\|\eta\|_2^2+C\sqrt L)
\]
is bounded by an absolute constant in this range, after increasing the
fixed lower threshold for \(L\).  Thus the damped family has norm at
most \(L^C S^{1/p}\) at every vertex, uniformly in all imaginary
parts.

Here are details that justify applying scalar three-lines to this
family.  Test against any element \(b\) of the dual unit ball
\(L^{p'}(Z;\mathcal H)\).  One may first restrict the pairing to a
bounded set on which every nonempty mass used in the finite family is
bounded below by a positive constant.  Such sets exhaust \(Z\), and
the restriction does not increase the norm of \(b\).  On each such
set, integration defines an entire scalar function; logarithms of the
finitely many masses are bounded there.  The crude bound above and the
damping make this scalar function bounded on every interpolating
strip.  If two real vectors \(a_0,a_1\) have established uniform
bounds for all imaginary parts, apply the bounded scalar three-lines
lemma \cite[Lemma~3.3, p.~57]{Hajlasz2012} to
\[
 z\longmapsto
 \int_Z
  \mathcal D\big((1-z)a_0+za_1+i\eta\big)
  \big\langle
   \mathcal F_v\big((1-z)a_0+za_1+i\eta,\cdot\big),b
  \big\rangle\,\dd x,
\]
with the indicated preliminary restriction understood.  The boundary
lines have the required bounds because their imaginary vectors are
still arbitrary.  The real parts throughout the strip remain in the
simplex.  Starting with its vertices and successively taking convex
combinations proves the same bound at their barycenter, with no
factor depending on \(K\).  Exhaustion of \(Z\), followed by duality,
removes the preliminary restriction.  Finally
\(\mathcal D(\zeta^{\mathrm c})=\exp(\mu^2/K)\) is an absolute
factor.  This proves the lemma.
\end{proof}

\subsection{The central weight and the exponent budget}

The central coefficient identity \eqref{eq:src-30} gives weights
\(\kappa_s\in(0,1]\).  Define
\[
 \mathcal W_\vartheta
       =\sum_s\kappa_s q_{T(s)}\phi_s.
\]
By that identity, its portion in root \(v\) is
\(n_v^{\mu/K}A_v(\zeta^{\mathrm c})\).

Put \(\gamma=1-\alpha/p+\mu/K\), which is positive for the values
of \(L\) under consideration.  Interpolation and H\"older's
inequality give, for each root,
\begin{align*}
 \int_Z\big\|n_v^{\mu/K}A_v(\zeta^{\mathrm c})\big\|
 &\le
  \|n_v^{\alpha/p-1}A_v(\zeta^{\mathrm c})\|_{L^p(Z)}
  \left(\int_Z n_v^{\gamma p'}\right)^{1/p'}\\
 &\le L^C S^{1/p}|F|^{1-1/p}(NL^C)^\gamma.
\end{align*}
Only an absolute number of roots meet the bounded set of top slopes.
Since \(S\lesssim L^C|E|=L^C N^2|F|\), summing over them proves
\begin{equation}\label{assembly:central-estimate}
 \int_{Z_\vartheta}\|\mathcal W_\vartheta\|
       \le L^C|F|N^{\,1+(2-\alpha)/p+\mu/K}.
\end{equation}

We record explicitly why this improves by every fixed negative power
of \(L\).  The net decrease of the exponent of \(N\) is
\[
 \delta_L
   =\frac{\alpha-2}{p}-\frac\mu K
   =\frac1p\left(\alpha-2-\frac{\alpha-1}{K}\right).
\]
The separated-count bound is a fixed power of \(L\), so the
definition of \(\alpha\) gives \(\alpha-2\ge c_1/\log L\) with a
fixed \(c_1>0\).  Also \(\alpha-1\le2\) and
\(K\ge(\log L)^2\).  For sufficiently large \(L\), therefore,
\[
 \delta_L\ge\frac{c_1}{2\sqrt L\log L},
 \qquad
 \delta_L\log N\ge\frac{c_1\sqrt L}{4\log L}.
\]
In particular, for every fixed \(Q>0\), after increasing the absolute
lower threshold for \(L\), Equation~\eqref{assembly:central-estimate}
implies, uniformly in \(\vartheta\),
\begin{equation}\label{assembly:central-power-gain}
 \int_{Z_\vartheta}\|\mathcal W_\vartheta\|
           \le L^{-Q}N|F|.
\end{equation}
This conclusion absorbs any fixed polynomial losses in the preceding
estimates; none of their exponents is changed here.

\subsection{Boundary averaging with the eligibility condition}

We now restore the dependence on the uniformly translated angular
grid.  The next estimate is pointwise on the fixed set \(Z_0\).

\begin{lemma}[Expected boundary loss]\label{assembly:boundary-loss}
Let \(\epsilon=4G^{-1/4}\).  For \(z\in Z_0\),
\begin{equation}\label{assembly:expected-error}
 \E_\vartheta\|P_q(z)-\mathcal W_\vartheta(z)\|
                    \le L^C\epsilon n(z).
\end{equation}
Here the expectation is normalized Lebesgue measure for the common
translation modulo \(d(0)\); independent translations at different
levels are neither assumed nor needed.
\end{lemma}

\begin{proof}
Fix \(z\in Z_0\) and a top \(T\), and suppress \(z\) from the
masses.  For \(i<k_T\) let
\[
 r_i(\vartheta)
       =\log\frac{m_{v_{i+1}(T)}}{n_{v_i(T)}}\ge0,
 \qquad
 M_i(b)=\sum_{\substack{U:\,|s_U-s_T|\le b\,d(i)\\
                              d_U\le d(i)/B^2}}h_U.
\]
Unlike \(r_i\), the latter mass does not depend on the grid
translation.  Each mass in this display is at least \(h_T\) when
\(i<k_T\), and is at most \(n\).

Let \(b_i(\vartheta)\in[0,1/2]\) be the distance from \(s_T\) to
the nearer boundary of its level-\(i\) interval, divided by \(d(i)\).
An extra top counted in
\(m_{v_{i+1}(T)}-n_{v_i(T)}\) has slope within
\(\epsilon d(i)\) of \(s_T\).  Indeed the neighbor distance in
Equation~\eqref{eq:src-14} is \(G^{-1/4}d(i)\), and the two
level-\((i+1)\) interval lengths add at most \(2d(i)/G\).
Such a top is eligible at level \(i+1\), hence also at level \(i\).
It lies outside the immediate parent.  It follows that \(r_i=0\)
if \(b_i>\epsilon\).  If \(0<b_i\le\epsilon\), the immediate
parent contains every top counted by \(M_i(b_i)\), whereas its extra
mass is disjoint from these inner tops and is at most
\(M_i(\epsilon)-M_i(b_i)\).  Consequently
\begin{equation}\label{assembly:one-boundary-loss}
 r_i(\vartheta)
       \le\log\frac{M_i(\epsilon)}{M_i(b_i(\vartheta))}.
\end{equation}
Endpoint coincidences can be omitted: for the finite set of slopes
and levels they form a null set of translations.

For a fixed number \(0<b<\epsilon\), set
\(r=\lceil\log_G(\epsilon/b)\rceil\).  The inclusion
\begin{equation}\label{assembly:mass-inclusion}
 M_{i+r}(\epsilon)\le M_i(b)
\end{equation}
uses both conditions in the definition of \(M_i\): the radius
satisfies \(\epsilon d(i+r)\le b d(i)\), and the precision threshold
\(d(i+r)/B^2\) is no larger than \(d(i)/B^2\).
To handle the end of the path without losing a term, define
\[
 a_i=\begin{cases}
       \log(M_i(\epsilon)/h_T),&0\le i<k_T,\\
       0,&i\ge k_T.
      \end{cases}
\]
If \(i+r<k_T\), Equation~\eqref{assembly:mass-inclusion} applies.
If \(i+r\ge k_T\), use \(M_i(b)\ge h_T\) instead.  In both cases
\[
 \log\frac{M_i(\epsilon)}{M_i(b)}\le a_i-a_{i+r}.
\]
Summing and cancelling gives
\begin{equation}\label{assembly:boundary-telescope}
 \sum_{i<k_T}\log\frac{M_i(\epsilon)}{M_i(b)}
 \le\sum_{0\le i<\min(r,k_T)}a_i
 \le r\log\frac n{h_T}.
\end{equation}
This argument also covers \(k_T=0\), with empty sums.

Because \(d(0)/d(i)=G^i\) is an integer, the common uniform
translation is uniform modulo every \(d(i)\).  Thus each
\(b_i\) has density \(2\) on \([0,1/2]\).  Linearity of
expectation, Equation~\eqref{assembly:one-boundary-loss}, and
Equation~\eqref{assembly:boundary-telescope} imply
\begin{align}
 \E_\vartheta\sum_{i<k_T}r_i(\vartheta)
 &\le2\int_0^\epsilon
       \sum_{i<k_T}\log\frac{M_i(\epsilon)}{M_i(b)}\,\dd b\notag\\
 &\le2\log(n/h_T)
          \int_0^\epsilon\lceil\log_G(\epsilon/b)\rceil\,\dd b
 \le C\epsilon\log(n/h_T).
 \label{assembly:total-boundary-loss}
\end{align}
For completeness the integral in the last line equals
\(\epsilon/(1-G^{-1})\): express the ceiling as a sum of indicators
of \(b<\epsilon G^{-j}\), \(j\ge0\), and integrate.  The dependence
of the different \(b_i\)'s is absent from this calculation.

On the fixed top \(T\), the value of \(\kappa_s\) depends only on
the stopping level \(k_s\).  As packet length increases, \(k_s\)
increases; Equation~\eqref{eq:src-30} shows that these values are
nonincreasing.  Their deficits from one are nondecreasing and satisfy
\[
 0\le1-\kappa_s
       \le\frac\mu K\sum_{i<k_T}r_i(\vartheta),
\]
using \(k_s\le k_T\) and \(1-e^{-x}\le x\).
Group the packets by stopping level.  Every interval of these levels
is an interval of packet lengths, so Equation~\eqref{eq:src-11}
bounds its unweighted sum by \(L^C h_T^8\).  Summation by parts
against the monotone deficits, and \(|q_T|\le1\), therefore give
\[
 \left\|\sum_{s:T(s)=T}(1-\kappa_s)q_T\phi_s\right\|
  \le L^C h_T^8\frac\mu K\sum_{i<k_T}r_i(\vartheta).
\]
For example, write the deficits as sums of their nonnegative
increments and use the bound on the corresponding suffix sums; the
sum of these increments is their largest value.  Taking expectations,
using \(\mu/K\le1\), and summing over tops proves
\[
 \E_\vartheta\|P_q-\mathcal W_\vartheta\|
       \le L^C\epsilon\sum_T h_T^8\log(n/h_T).
\]
The weights have a fixed absolute upper bound.  For every such weight,
\[
 h_T^8\log(n/h_T)
       \le C h_T(1+\log_+ n),
\]
because \(h^7\log_+(1/h)\) is bounded on a fixed bounded interval.
Equation~\eqref{eq:src-12} gives \(1+\log_+n\lesssim L\) on
\(Z_0\).  This proves Equation~\eqref{assembly:expected-error}.
\end{proof}

\subsection{Recovering the original adjoint sum}

\begin{theorem}[Improvement for the popular model]
\label{assembly:improved-model}
The popular adjoint sum satisfies Equation~\eqref{eq:src-7}; namely,
\[
 \int_F\|P\|\lesssim L^{-20}N|F|.
\]
\end{theorem}

\begin{proof}
We make the exceptional-set comparison explicit because the good set
depends on the grid.  Write
\(\|P\|_2\le L^{C_P}\sqrt{|E|}\), as supplied by
Equation~\eqref{eq:src-4}.  Proposition~\ref{conc:mask-removal} and
Proposition~\ref{count:simultaneous} allow the fixed polynomial thresholds
to be chosen so that
\begin{equation}\label{assembly:fixed-exceptions}
 \int_{F\setminus Z_0}\|P\|\le L^{-30}N|F|,
 \qquad
 \int_{Z_0}\|P-P_q\|\le L^{-30}N|F|.
\end{equation}
Indeed a discarded area at most
\(L^{-2C_P-60}|F|\) has the stated first bound by Cauchy--Schwarz.
The exponentially small concentration and count exceptions meet this
area bound for sufficiently large \(L\).

For each fixed translation the good-path construction gives
\[
 |Z_0\setminus Z_\vartheta|\le S e^{-9L}.
\]
This is Markov's inequality applied to the weighted bad-path mass:
its integral is at most \(S e^{-10L}\), and the deletion threshold
is \(e^{-L}\).  Since \(S\lesssim L^C N^2|F|\) and
\(N^2\le e^{4L}\),
\begin{equation}\label{assembly:moving-exceptions}
 \int_{Z_0\setminus Z_\vartheta}\|P\|
       \le L^{C_P}\sqrt{|E|}\,(S e^{-9L})^{1/2}
       \le L^C e^{-5L/2}N|F|
       \le L^{-30}N|F|.
\end{equation}
These are bounds for each translation; no union over translations is
taken.

For every \(\vartheta\), the triangle inequality gives the following
decomposition, whose left side is independent of \(\vartheta\):
\begin{align}
 \int_F\|P\|
 &\le \int_{F\setminus Z_\vartheta}\|P\|
      +\int_{Z_\vartheta}\|P-P_q\|
      +\int_{Z_\vartheta}\|P_q-\mathcal W_\vartheta\|
      +\int_{Z_\vartheta}\|\mathcal W_\vartheta\|.
 \label{assembly:original-decomposition}
\end{align}
The first two terms are controlled by
Equations~\eqref{assembly:fixed-exceptions} and
\eqref{assembly:moving-exceptions}.  The final term is controlled by
Equation~\eqref{assembly:central-power-gain} with \(Q=30\).
Average the remaining term over translations.  It is nonnegative and
\(Z_\vartheta\subset Z_0\), so Tonelli's theorem and
Lemma~\ref{assembly:boundary-loss} yield
\begin{align*}
 \E_\vartheta\int_{Z_\vartheta}
       \|P_q-\mathcal W_\vartheta\|
 &\le\int_{Z_0}\E_\vartheta
       \|P_q-\mathcal W_\vartheta\|\\
 &\le L^C\epsilon\int_{Z_0}n
 \le L^C\epsilon N|F|
 \le L^{-30}N|F|.
\end{align*}
The penultimate inequality uses Equation~\eqref{eq:src-12}; the last
uses \(\epsilon=4G^{-1/4}\) and \(\log G\asymp L^{.7}\).
All functions being averaged are measurable: for finitely many tops,
packets, and levels, the angular assignments change only at finitely
many translation values, and the good sets were defined by measurable
quantities.  Thus this averaging is legitimate even though
\(Z_\vartheta\) and \(\mathcal W_\vartheta\) use the same
translation.  Averaging Equation~\eqref{assembly:original-decomposition}
now proves the asserted estimate.  In particular, on every discarded
set we have estimated the original sum \(P\).
\end{proof}

\subsection{Choice of parameters and completion of the reductions}

We collect the order of the absolute choices.  First fix the packet
adaptation, separation, and overlap constants and the fixed geometric
enlargements used with \Cref{red:kakeya}.  Fix the local slope bound
\(C_{\rm Lip}\), strip length \(b_0\), and outer length \(a=a_*\)
as in \eqref{red:local-operator}, with the required room for those
enlargements.  Choose the low-density exponent \(C_2\) by
\eqref{tile:eq-improvement-loss}, before constructing the popular forest.

Next fix the decay exponent in the top-count weight and the concentration
orders that determine \(K_1\).  Choose the fixed mask-removal and
exceptional-area powers, then choose \(A\) using the count at dilation
\(K_1\), whose constants do not depend on \(A\).  The count at the larger
dilation \(K_0=L^{A+2}\) is used only after this comparison.  Fix the threshold
in \eqref{eq:src-12}, the separated-count exponent defining \(D\), and
the good-path thresholds.  These choices determine the polynomial powers
in the motion, prediction, and residual estimates.

Now choose \(M=L^{A_M}\) as in \eqref{res:choice-M}.  For large \(L\),
\(\mu<1/12\), so this fixed exponent controls both
\[
 pL^C M^{-1/2+\mu}
 \quad\text{and}\quad
 pL^C M^{-1/3+\mu}
\]
in the marked and unary residual estimates.  The good-path thresholds
were fixed before \(A_M\), and their required proportion remains valid:
\[
 e^{-L}L^{Cp}(1+M)^{\alpha-1}
       =\exp\big(-L+O(\sqrt L\log L)+O(\log L)\big)<\tfrac12
\]
for sufficiently large \(L\).

Finally choose \(L_0\) after the outer length and all fixed exponents.
It ensures \(\tau<a\), the geometric smallness requirements, and all
the inequalities used above, because
\[
 \log B=O(L^{.7}(\log L)^2)=o(L^{.9}),\qquad
 \log(N^\mu)=O(\sqrt L)=o(\log G),
\]
and, for every fixed exponent \(C\),
\[
 \log(m^C\tau)=O(L^{.9})-\Theta(L^{.97})\longrightarrow-\infty.
\]
It also enforces the saving in \eqref{assembly:central-power-gain}.
All choices are independent of the depth, tile count, and dimension.

\begin{theorem}[The three model estimates]
\label{assembly:all-model-estimates}
For the packet forms in \eqref{eq:src-3}, let \(u\) be a bounded
measurable real slope, let \(F,E\subset\R^2\) be measurable sets of
positive finite measure, and suppose
\(\|f\|_{\mathcal H}\le\one_F\) and
\(\|g\|_{\mathcal H}\le\one_E\) pointwise.  The following estimates
hold uniformly over every finite subcollection.
\begin{enumerate}[label=\textup{(\roman*)}]
\item For \(m\ge1\),
\[
 |\Lambda(f,g)|\lesssim(1+\log(2m))^C\sqrt{|F||E|}.
\]
\item Suppose \(L\ge L_0\), with \(\tau,m,W\) chosen by
\eqref{red:small-time-parameters}, and
\(L/2\le\log\sqrt{|E|/|F|}\le2L\).  If
\(\abs u\le1\), \(\Lip(u)\le C_{\rm Lip}\), and
\(F,E\subset I\times\R\) for an interval \(I\) of length at most
\(b_0\), then
\[
 |\Lambda(f,g)|\lesssim L^{-20}\sqrt{|F||E|}.
\]
\item For the scalar model with \(m=1\) and \(|E|\le|F|\), there is an absolute
\(c>0\) such that
\[
 |\Lambda(f,g)|\lesssim
       \sqrt{|F||E|}\left(\frac{|E|}{|F|}\right)^c.
\]
\end{enumerate}
\end{theorem}

\begin{proof}
The first and third assertions are precisely the ordinary estimates
proved with \eqref{eq:src-4}.  For the second, the fixed deletion
of tiles with \(D_E(s)\le L^{-C_2}\) in \Cref{forest:section}
contributes at most \(CL^{-40}\sqrt{|F||E|}\), by
\eqref{tile:eq-improvement-loss}.  Assign the remaining tiles to the popular tops
as already constructed.  Their adjoint sum is \(P\), and the
coefficient convention in Equation~\eqref{eq:src-3} gives
\[
 |\Lambda_{\mathrm{popular}}(f,g)|
       \le\int_F\|P\|.
\]
Theorem~\ref{assembly:improved-model} bounds this by
\(L^{-20}N|F|=L^{-20}\sqrt{|F||E|}\).  Adding the deleted
contribution proves the second assertion.  All preceding arguments
apply to arbitrary finite tile sets and thus give the claimed
uniformity.
\end{proof}

\begin{proof}[Completion of the proofs of \Cref{thm:band-restricted,thm:main}]
Theorem~\ref{assembly:all-model-estimates} supplies exactly the three
hypotheses of Proposition~\ref{red:model-suffices}.  That proposition
therefore proves the band restricted estimate
\eqref{eq:src-1}.  The dependence here is forward: the
reductions established the implication from the model estimates,
which have now been proved without using the band restricted theorem.

All estimates were obtained first for finite packet collections and
positive inner truncations, with constants independent of both.
The summable smooth expansions and translated-lattice averages in the
packet reduction consequently pass to their limits as specified in
Proposition~\ref{red:model-suffices}.  The amplitude summation,
vertical Littlewood--Paley reassembly, and spatial localization proved
in the reductions give a uniform strong \(L^2\) estimate for the
original fixed-field transform with a positive inner cutoff.  For a
Schwartz input its symmetric principal value exists pointwise; Fatou's
lemma gives the same strong bound for the principal-value operator, which
extends by density.  For the already chosen outer length \(a=a_*\),
Chebyshev's inequality gives the weak \((2,2)\) assertion
of Theorem~\ref{thm:main}.  The constants are independent of the
vector field, the input, and the level in that inequality. The dilation
calculation following \Cref{thm:main} gives the outer length
\(a_*/\Lip(v)\) for positive Lipschitz seminorm; constant directions
have already been handled by the one-dimensional theorem.
\end{proof}
